% =====================================================================
\chapter{Use case của hệ thống}\label{chap:usecase}
% =====================================================================

\section{Tác nhân}
\begin{table}[H]
\centering\small
\caption{Các tác nhân của hệ thống}\label{tab:actors}
\begin{tabularx}{\linewidth}{@{}lX@{}}
\toprule
\textbf{Tác nhân} & \textbf{Mô tả} \\ \midrule
Người lái xe & Người gửi/lấy xe; dùng thẻ RFID và nhìn màn hình TFT tại bãi. \\
Người vận hành & Bảo vệ hoặc quản lý bãi; dùng trang web trên điện thoại/máy tính để giám sát và điều khiển thanh chắn. \\
Sensor Client & Tác nhân thiết bị: bo ESP8266 đọc cảm biến chỗ đỗ, lửa, nước và gửi dữ liệu lên máy chủ. \\
RFID Client & Tác nhân thiết bị: bo ESP8266 đọc thẻ ở hai cổng và cảm biến nhiệt độ; đồng thời phát mạng Wi-Fi. \\
Đồng hồ hệ thống & Tác nhân thời gian: kích hoạt việc tự đóng thanh chắn, hết hạn hóa đơn/thông báo. \\
\bottomrule
\end{tabularx}
\end{table}

\section{Sơ đồ use case}
Hình~\ref{fig:usecase} tổng hợp 10 use case và quan hệ với các tác nhân. Hai tác nhân
thiết bị (bên phải) là nguồn dữ liệu; hai tác nhân con người (bên trái) là người
hưởng lợi.

\begin{figure}[H]
\centering
\mmd[1][0.62]{usecase}
\caption{Sơ đồ use case tổng quát (đường nét đứt: sự kiện theo thời gian — tự đóng thanh chắn, hết hạn thông báo)}\label{fig:usecase}
\end{figure}

\section{Bảng tổng hợp use case}
\begin{table}[H]
\centering\small
\caption{Tổng hợp 10 use case, nơi hiện thực và bộ kiểm thử tương ứng}\label{tab:uc-summary}
\begin{tabularx}{\linewidth}{@{}l>{\raggedright\arraybackslash}p{4.2cm}>{\raggedright\arraybackslash}X>{\raggedright\arraybackslash}p{3.1cm}@{}}
\toprule
\textbf{Mã} & \textbf{Tên use case} & \textbf{Module chính} & \textbf{Kiểm thử} \\ \midrule
UC01 & Theo dõi chỗ trống tại bãi & \file{lcd\_display.c} & \file{sensor\_client} \\
UC02 & Giám sát bãi qua web & \file{http\_server.c}, \file{index.html} & S2, S12 \\
UC03 & Gửi xe bằng thẻ RFID & \file{rfid\_session.c}, \file{servo.c} & \file{rfid\_session}, \file{rfid\_client} \\
UC04 & Lấy xe và thanh toán phí & \file{rfid\_session.c} & \file{rfid\_session} \\
UC05 & Từ chối quẹt sai cổng / bãi đầy & \file{rfid\_session.c}, \file{lcd\_display.c} & \file{rfid\_session} \\
UC06 & Cảnh báo cháy và sơ tán & \file{udp\_server.c}, \file{servo.c}, \file{buzzer.c} & \file{sensor\_client}, S9 \\
UC07 & Giám sát mực nước & \file{SensorClient.ino} & \file{sensor\_client} \\
UC08 & Giám sát nhiệt độ & \file{RFIDClient.ino}, \file{parking\_state.c} & \file{rfid\_client}, \file{rfid\_session} \\
UC09 & Điều khiển thanh chắn thủ công & \file{http\_server.c}, \file{servo.c} & S12, S13 \\
UC10 & Xem nhật ký vào/ra và giá & \file{rfid\_session.c}, \file{http\_server.c} & \file{rfid\_session} \\
\bottomrule
\end{tabularx}
\end{table}
\noindent Tên tệp kiểm thử như \file{rfid\_session} chỉ bộ kiểm thử tự động
(Chương~\ref{chap:kiemthu}); S2, S9,\ldots{} là kịch bản kiểm thử trên phần cứng
(Bảng~\ref{tab:systest}).

\section{Đặc tả chi tiết}

\begin{usecase}{UC01}{Theo dõi chỗ trống tại bãi}
\item[Tác nhân] Người lái xe (chính); Sensor Client (nguồn dữ liệu).
\item[Mô tả] Người lái nhìn màn hình TFT ở lối vào để biết số chỗ trống và vị trí từng chỗ trống.
\item[Tiền điều kiện] Sensor Client và máy chủ đã kết nối Wi-Fi \code{ESP32-Parking}.
\item[Luồng chính] \begin{enumerate}[nosep,leftmargin=1.3em]
  \item Xe đi vào/ra một ô đỗ, cảm biến IR của ô đó đổi mức logic.
  \item Sensor Client lọc nhiễu: chỉ chấp nhận mức mới khi giữ ổn định 50\,ms.
  \item Ngay khi trạng thái đổi, Sensor Client gửi bản tin \code{parking\_update} (UDP).
  \item Máy chủ kiểm tra bản tin, cập nhật trạng thái chung và đánh thức tác vụ màn hình.
  \item Màn hình vẽ lại: số chỗ trống \code{FREE SPACES x/5} và 5 ô S1–S5 (xanh = trống, đỏ = có xe).
\end{enumerate}
\item[Ngoại lệ] 3a. Mất gói UDP: bản tin nhắc lại sau tối đa 1\,s sửa lại trạng thái.
  4a. Bản tin sai định dạng: bị bỏ qua, trạng thái cũ được giữ.
\item[Hậu điều kiện] Màn hình phản ánh đúng trạng thái 5 chỗ đỗ.
\item[Thông số] Debounce 50\,ms; nhắc lại 1\,s; độ trễ hiển thị $\approx$ 0,1\,s.
\end{usecase}

\begin{usecase}{UC02}{Giám sát bãi qua trang web}
\item[Tác nhân] Người vận hành.
\item[Mô tả] Xem trực quan sơ đồ bãi, số chỗ trống, cảm biến và trạng thái thanh chắn trên điện thoại.
\item[Tiền điều kiện] Thiết bị của người vận hành kết nối Wi-Fi \code{ESP32-Parking}.
\item[Luồng chính] \begin{enumerate}[nosep,leftmargin=1.3em]
  \item Người vận hành mở \code{http://parking.local}.
  \item Thiết bị phân giải tên bằng mDNS và nhận địa chỉ IP hiện tại của máy chủ.
  \item Máy chủ trả trang web nhúng trong firmware.
  \item Trang web gọi \code{GET /api/parking} mỗi 1\,s và vẽ sơ đồ 5 ô đỗ (xe màu đỏ, ô trống chữ P màu xanh).
\end{enumerate}
\item[Ngoại lệ] 2a. Thiết bị không hỗ trợ \code{.local}: dùng địa chỉ IP in trên log máy chủ.
  4a. Máy chủ mất kết nối: trang hiển thị \code{ESP32 / OFFLINE}.
\item[Hậu điều kiện] Người vận hành thấy trạng thái bãi với độ trễ $\le$ 1\,s.
\end{usecase}

\begin{usecase}{UC03}{Gửi xe bằng thẻ RFID}
\item[Tác nhân] Người lái xe; RFID Client; Đồng hồ hệ thống.
\item[Mô tả] Người lái quẹt thẻ ở cổng vào để được mở thanh chắn và bắt đầu tính giờ.
\item[Tiền điều kiện] Thẻ chưa có phiên gửi xe; bãi còn phiên trống (tối đa 5).
\item[Luồng chính] \begin{enumerate}[nosep,leftmargin=1.3em]
  \item Người lái đưa thẻ vào đầu đọc cổng vào.
  \item RFID Client đọc UID, chuyển thành chuỗi hex in hoa, gửi \code{rfid\_event} với \code{"gate":"in"}.
  \item Máy chủ tạo phiên mới, lưu thời điểm vào, ghi nhật ký \code{IN}.
  \item Máy chủ mở thanh chắn cổng vào (90°) và hẹn giờ đóng sau 3\,s.
  \item Đồng hồ hết 3\,s, thanh chắn tự đóng.
\end{enumerate}
\item[Ngoại lệ] 1a. Thẻ để yên trên đầu đọc: chỉ gửi một lần trong 2\,s.
  3a. Bãi đầy hoặc thẻ đã ở trong bãi: chuyển sang UC05.
  5a. Người vận hành bấm đóng/mở trên web trước khi hết giờ: lệnh tay thắng, hủy hẹn giờ.
\item[Hậu điều kiện] Thẻ có phiên hoạt động; thanh chắn đóng lại.
\item[Thông số] Thời gian mở 3\,s; chống đọc lặp 2\,s; tối đa 5 phiên.
\end{usecase}

\begin{usecase}{UC04}{Lấy xe và thanh toán phí}
\item[Tác nhân] Người lái xe; RFID Client; Đồng hồ hệ thống.
\item[Mô tả] Người lái quẹt cùng thẻ ở cổng ra; hệ thống tính phí, hiển thị hóa đơn và mở cổng ra.
\item[Tiền điều kiện] Thẻ đang có phiên gửi xe.
\item[Luồng chính] \begin{enumerate}[nosep,leftmargin=1.3em]
  \item Người lái quẹt thẻ ở đầu đọc cổng ra (\code{"gate":"out"}).
  \item Máy chủ tính thời gian gửi $t$ và phí $\max(1,\lceil t/3600\rceil)\times5000$\,VNĐ.
  \item Máy chủ đóng phiên, ghi nhật ký \code{OUT} kèm thời gian và phí.
  \item Màn hình hiển thị hóa đơn: UID, \code{PARK TIME}, \code{FEE VND} trong 60\,s.
  \item Thanh chắn cổng ra mở và tự đóng sau 3\,s.
\end{enumerate}
\item[Ngoại lệ] 1a. Thẻ chưa vào bãi: chuyển sang UC05.
\item[Hậu điều kiện] Phiên được giải phóng; nhật ký có bản ghi \code{OUT}.
\item[Thông số] 5.000\,VNĐ/giờ, làm tròn lên, tối thiểu 1 giờ; hóa đơn hiển thị 60\,s.
\end{usecase}

\begin{usecase}{UC05}{Từ chối quẹt sai cổng / bãi đầy}
\item[Tác nhân] Người lái xe; RFID Client; Đồng hồ hệ thống.
\item[Mô tả] Hệ thống từ chối các lần quẹt không hợp lệ và hướng dẫn người lái.
\item[Tiền điều kiện] Một trong ba trường hợp: (a) thẻ đã trong bãi nhưng quẹt ở cổng vào;
  (b) thẻ chưa vào nhưng quẹt ở cổng ra; (c) đã đủ 5 phiên mà có thẻ mới ở cổng vào.
\item[Luồng chính] \begin{enumerate}[nosep,leftmargin=1.3em]
  \item Máy chủ nhận \code{rfid\_event} và xác định kết quả \code{INVALID\_GATE} hoặc \code{FULL}.
  \item Thanh chắn \emph{không} mở.
  \item Màn hình hiển thị thông báo đỏ trong 5\,s: \code{WRONG GATE} kèm gợi ý
        \code{ALREADY INSIDE - USE EXIT GATE} / \code{NOT CHECKED IN - USE ENTRY GATE},
        hoặc \code{PARKING FULL}.
  \item Nhật ký ghi \code{REJECTED · WRONG GATE} hoặc \code{REJECTED · FULL}.
\end{enumerate}
\item[Ngoại lệ] 3a. Đang hiển thị hóa đơn: thông báo hiện trước, hết 5\,s hóa đơn hiện lại.
  3b. Ngay sau đó có lần quẹt hợp lệ: thông báo cũ bị xóa ngay.
\item[Hậu điều kiện] Trạng thái phiên không đổi.
\end{usecase}

\begin{usecase}{UC06}{Cảnh báo cháy và sơ tán}
\item[Tác nhân] Sensor Client (cảm biến lửa); Người lái xe; Người vận hành.
\item[Mô tả] Khi phát hiện lửa: báo động tại chỗ, trên web và mở thanh chắn để sơ tán.
\item[Tiền điều kiện] Sensor Client đang gửi dữ liệu.
\item[Luồng chính] \begin{enumerate}[nosep,leftmargin=1.3em]
  \item Cảm biến lửa phát hiện ngọn lửa ổn định quá 50\,ms; Sensor Client gửi ngay \code{"fire":true}.
  \item Máy chủ bật còi (GPIO9 mức thấp).
  \item Máy chủ chuyển \emph{chế độ cháy}: mở cả hai thanh chắn và giữ mở.
  \item Màn hình chuyển sang màn hình \code{DANGER / FIRE DETECTED / EVACUATE NOW} có ngọn lửa động.
  \item Trang web hiện băng rôn đỏ \code{FIRE DETECTED}.
\end{enumerate}
\item[Ngoại lệ] 3a. Người vận hành bấm đóng thanh chắn: bị từ chối (HTTP 500) trong chế độ cháy.
  6a. Hết lửa (\code{"fire":false}): tắt còi, đóng cả hai thanh chắn, màn hình trở lại bình thường.
\item[Hậu điều kiện] Lối ra vào thông thoáng trong suốt thời gian có lửa.
\end{usecase}

\begin{usecase}{UC07}{Giám sát mực nước}
\item[Tác nhân] Sensor Client; Người vận hành.
\item[Mô tả] Theo dõi mực nước ngập tại bãi (mô hình tầng hầm).
\item[Luồng chính] \begin{enumerate}[nosep,leftmargin=1.3em]
  \item Sensor Client đọc ADC chân A0 mỗi 250\,ms.
  \item Quy đổi tuyến tính ADC 20–700 thành 0–40\,mm (giới hạn trên 40\,mm).
  \item Giá trị được gửi trong \code{parking\_update}; máy chủ kiểm tra dải 0–10.000\,mm.
  \item Màn hình (\code{WATER MM}) và web (\code{WATER LEVEL}) hiển thị giá trị.
\end{enumerate}
\item[Ngoại lệ] 3a. Giá trị âm hoặc vượt dải: bản tin bị từ chối.
\item[Hậu điều kiện] Mực nước hiển thị với độ phân giải 0,1\,mm.
\end{usecase}

\begin{usecase}{UC08}{Giám sát nhiệt độ}
\item[Tác nhân] RFID Client (DHT11); Người vận hành.
\item[Mô tả] Theo dõi nhiệt độ môi trường tại cổng.
\item[Luồng chính] \begin{enumerate}[nosep,leftmargin=1.3em]
  \item RFID Client đọc DHT11 mỗi 5\,s, giữ giá trị hợp lệ gần nhất.
  \item Giá trị được gửi kèm mỗi \code{rfid\_event} (\code{temperature\_c}).
  \item Máy chủ kiểm tra dải $-40$ đến $125$\,°C, cập nhật trạng thái chung.
  \item Màn hình (\code{TEMP C}) và web (\code{TEMPERATURE}) hiển thị.
\end{enumerate}
\item[Ngoại lệ] 1a. DHT11 đọc lỗi (NaN): giữ giá trị cũ, không gửi số sai.
\item[Hậu điều kiện] Nhiệt độ cập nhật mỗi lần có xe quẹt thẻ.
\end{usecase}

\begin{usecase}{UC09}{Điều khiển thanh chắn thủ công}
\item[Tác nhân] Người vận hành.
\item[Mô tả] Mở/đóng từng thanh chắn hoặc cả hai (khẩn cấp) từ trang web.
\item[Luồng chính] \begin{enumerate}[nosep,leftmargin=1.3em]
  \item Người vận hành bấm \code{OPEN}/\code{CLOSE} của cổng vào, cổng ra, hoặc
        \code{EMERGENCY OPEN/CLOSE ALL}.
  \item Trang web gửi yêu cầu \code{POST} tới \code{/api/barrier/<cổng>/<lệnh>} (cổng: \code{in}, \code{out}, \code{all}; lệnh: \code{open}, \code{close}); nút bị khóa đến khi có phản hồi.
  \item Máy chủ hủy hẹn giờ tự đóng (nếu có) rồi đặt góc servo.
  \item Trang web cập nhật trạng thái \code{OPEN}/\code{CLOSED}.
\end{enumerate}
\item[Ngoại lệ] 3a. Đang có cháy mà yêu cầu đóng: từ chối, hiển thị lỗi.
\item[Hậu điều kiện] Thanh chắn ở đúng trạng thái người vận hành chọn.
\end{usecase}

\begin{usecase}{UC10}{Xem nhật ký vào/ra và giá gửi xe}
\item[Tác nhân] Người vận hành.
\item[Mô tả] Biết ai đang trong bãi, đã gửi bao lâu, phí tạm tính, và 10 sự kiện gần nhất.
\item[Luồng chính] \begin{enumerate}[nosep,leftmargin=1.3em]
  \item Trang web gọi \code{GET /api/rfid} mỗi 2\,s.
  \item Máy chủ trả giá mỗi giờ, danh sách xe trong bãi (UID, giờ vào, phí tạm tính) và
        10 sự kiện mới nhất (vào, ra kèm thời lượng và phí, bị từ chối).
  \item Trình duyệt đổi thời gian ``giây kể từ khi khởi động'' thành giờ thực theo đồng hồ điện thoại.
\end{enumerate}
\item[Ngoại lệ] 2a. Máy chủ khởi động lại: nhật ký và phiên được làm mới (lưu trong RAM).
\item[Hậu điều kiện] Người vận hành đối soát được tiền thu và xe trong bãi.
\end{usecase}
